\begin{hebrew}
\begingroup
\let\raggedright\raggedleft
\chapter*{תקציר}
\endgroup
\begin{english}
\addcontentsline{toc}{chapter}{\HebrewTOC{תקציר}}
\end{english}

\raggedleft
\setlength{\parindent}{0pt}
\setlength{\parskip}{0.75\baselineskip}

בנתונים רפואיים מהעולם האמיתי אין לאומדים סיבתיים מפתח תשובות סיבתי נצפה לעובדות הנגד־עובדתיות, ואילו ערכות הערכה מבוקרות משיגות תשובה ידועה במחיר של פישוט או ייצור חלקים מתהליך יצירת הנתונים. \textenglish{CliniCause} בונה ערכות מבחן תצפיתיות משלימות לניתוח סיבתי, המשמרות מדידות לא־סדירות שנצפו במקור ביחידות טיפול נמרץ, דפוסי חסר, משתנים נצפים ותמותה, ומציגות במפורש את הייצוגים וההנחות שהוגדרו במחקר.

מודל שפה גדול, שעוגן בסכמות הנתונים ובספרות, סייע רק בשלב התכנון בהצעת מצבי־פרוקסי, משפחות כללים וגרפים סיבתיים; הוא לא נחשף לרשומות מטופלים ולא אמד ניגודים. הצעות שנבחרו במסגרת הפרויקט קודדו בכללים ובגרפים דטרמיניסטיים. שני המשאבים הנפרדים כוללים 26,845 רשומות ותשע חשיפות פרוקסי מנותחות ב־\textenglish{MIMIC-III}, ו־7,993 רשומות ועשר חשיפות ב־\textenglish{PhysioNet 2012}. בכל משאב נבחנו ארבעה מודלי סדרות זמן—\textenglish{STraTS}, ‏\textenglish{GRU}, ‏\textenglish{GRU-D} ו־\textenglish{TCN}—ושלושה אומדים סיבתיים: \textenglish{CausalForestDML} כאומד ראשי, \textenglish{LinearDML} כמשווה משני ו־\textenglish{CausalPFN} כאומד חקרני, ללא איגום בין המאגרים.

בחיזוי תוויות הפרוקסי היה טווח \textenglish{AUROC} במבחן \textenglish{0.867--0.918}; \textenglish{STraTS} הוביל במדדי \textenglish{MIMIC} הארכיוניים, ו־\textenglish{GRU-D} הוביל במדדי \textenglish{PhysioNet}. \textenglish{CausalForestDML} ו־\textenglish{LinearDML} הסכימו בכיוון הניגוד שנאמד במודל בכל 19 ההשוואות בין מאגר לחשיפה, ושלושת האומדים הסכימו ב־18 מתוך 19. הלם ב־\textenglish{PhysioNet} היה החריג: שני אומדני ה־\textenglish{DML} היו שליליים, ואומדן \textenglish{CausalPFN} היה חיובי קלות. השוואה הקלינית ההקשרית הצביעה הן על חפיפה רחבה והן על פערים מהותיים. הסכמת האומדים התקיימה לצד מגבלות התאמה ותמיכה, שינויים בעקבות הפרעת התוצאה ורגישות לערבול שהושמט.

לאף אחת ממערכות המבחן אין אמת סיבתית ידועה, ומצבי־הפרוקסי, הגרפים והניגודים שנאמדו במודל לא עברו תיקוף קליני. התוצאות אינן מבססות השפעות טיפול, המלצות טיפול או מוכנות להטמעה. אינטגרציה, חוזי נתונים מפורשים, שכבת ייצוג דטרמיניסטית, אבחונות ופרובננס הופכים את המשאבים לניתנים לבדיקה. תרומתם היא בהשלמת ערכות הערכה המבוקרות באמצעות חשיפת יציבות האומדים, אי־הסכמות, מגבלות תמיכה ורגישות בתנאים קליניים שנצפו במקור.
\end{hebrew}
